\documentclass[11pt,a4paper]{article}
\usepackage[utf8]{inputenc}
\usepackage[spanish]{babel}
\usepackage{geometry}
\geometry{top=2.5cm, bottom=2.5cm, left=2.5cm, right=2.5cm}
\usepackage{enumitem}

\title{\textbf{Minuta de Reunión: Programa de Formación Especializada}}
\author{}
\date{Febrero 2026}

\begin{document}
\maketitle
\vspace{-1.5cm}

\section{Metodología y Modalidad}
\begin{itemize}[noitemsep]
    \item \textbf{Modalidad híbrida:} Apartado teórico presencial/remoto y asesorías virtuales. Clases de laboratorio estrictamente presenciales.
    \item \textbf{Carga académica:} 140 horas de contenido, más trabajo extracurricular.
    \item \textbf{Enfoque práctico:} Las horas de asesoría se usarán para responder dudas sobre datos y problemas reales del ámbito laboral del estudiante.
    \item \textbf{Docencia:} Impartido por doctores especializados en franjas horarias definidas.
    \item \textbf{Progresión:} El diseño debe garantizar que el estudiante perciba un avance continuo en sus conocimientos.
\end{itemize}

\section{Evaluación y Perfil del Estudiante}
\begin{itemize}[noitemsep]
    \item \textbf{Público objetivo:} Personal específico de las fuerzas que requiere la formación para mejorar su rendimiento laboral. El curso no está abierto al público general.
    \item \textbf{Aprobación:} La nota aprobatoria se basará en la entrega de un proyecto final.
    \item \textbf{Certificación:} Se entregará certificado tras finalizar clases y aprobar.
\end{itemize}

\section{Seguridad y Manejo de Información}
\begin{itemize}[noitemsep]
    \item \textbf{Restricción de datos:} Prohibido el uso de bases de datos de interés nacional. Solo se utilizará información disponible en la red.
    \item \textbf{Protocolos:} Se deben respetar estrictamente los protocolos de seguridad de la información del Ejército.
\end{itemize}

\section{Contenido del Programa (Objetivo: Sistema Anti-UAV)}
El programa debe abarcar la diversidad de intereses de las fuerzas, dividiéndose en:
\begin{itemize}[noitemsep]
    \item \textbf{Módulo 1:} Naturaleza del espectro y arquitectura SDR.
    \item \textbf{Módulo 2:} FFT, análisis espectral, vigilancia espectral y prácticas.
    \item \textbf{Módulo 3:} Identificación de señales específicas y análisis estadístico.
    \item \textbf{Módulo 4:} Guerra electrónica, consideraciones regulatorias y \textit{jamming}.
    \item \textbf{Módulo 5:} Principios de radar, modelado matemático y detección de objetivos pequeños.
    \item \textbf{Módulo 6:} Inteligencia artificial, \textit{datasets} y Python para análisis.
    \item \textbf{Módulo 7:} Aprendizaje supervisado, \textit{Random Forest}, clasificación y curvas ROC.
    \item \textbf{Módulo 8:} \textit{Deep Learning}, detección automática de UAV y reducción de falsas alarmas.
    \item \textbf{Módulo 9:} Integración de conocimientos para el proyecto final.
\end{itemize}

\section{Recursos Necesarios}
\begin{itemize}[noitemsep]
    \item \textbf{Hardware:} Analizadores de espectro, SDR multibanda y generadores de señales.
\end{itemize}

\section{Acuerdos Institucionales y Tareas Pendientes}
\begin{itemize}[noitemsep]
    \item \textbf{Alianza Institucional:} Establecer un mecanismo menos formal para contar con asesoría de la Universidad Nacional enfocada en resolver problemas del Ejército.
    \item \textbf{Intercambio:} Explorar posibilidades de pasantías, entrega de material y visitas de personal. El Ejército hará la selección de su personal.
    \item \textbf{Gestión Administrativa:} 
    \begin{itemize}[noitemsep]
        \item Comunicar al decano con los directores de la escuela (alineando intereses comunes).
        \item Sacar el borrador final y hacer la convocatoria de estudiantes.
        \item Revisar a detalle la intensidad de las clases, la modalidad final y los costos asociados.
        \item Discutir la pertinencia de presentar esta propuesta directamente al Ministerio.
    \end{itemize}
\end{itemize}

\section{Próximas Reuniones}
\begin{itemize}[noitemsep]
    \item \textbf{Reunión de urgencia (Proyectos de alta prioridad):} Jueves, 19 de febrero.
    \item \textbf{Reunión de seguimiento general:} Lunes, 23 de febrero a las 3:00 PM.
\end{itemize}

\end{document}